\chapter{Toward a Simulation Framework}
\label{app:simulation}

\begin{quotation}
\noindent\textit{Synopsis.} This appendix outlines, without implementing,
how budget relays could be represented computationally and investigated
experimentally through simulation.
\end{quotation}


This appendix sketches, without implementing, how a budget relay
(\cref{def:budget-relay}) could be simulated computationally. It is
explicitly scoped as a direction for future computational work, not
as part of the monograph's formal content.

Each stage of a relay would be represented as a process holding a
declared triple $(C_i, I_i, \beta_i)$, applying
\cref{def:budget-discard} to a set of input distinctions and producing
an output rendering together with a loss ledger $L_i$ per
\cref{def:loss-ledger}. A downstream query interface would then test
\cref{thm:composition-failure}'s failure mode directly, by
constructing task instances $\tau^*$ outside the set of tasks any
stage was configured to anticipate, and checking whether the final
rendering's admissible continuation set $\admis{\tau^*}{A_n}$ contains
the correct continuation.

Such a framework would allow empirical (rather than purely
constructive) exploration of \cref{cor:no-free-lunch}'s conjectured
growth bound across relays of varying length, branching factor
(\cref{ch:fanout}), and coordination mechanism
(\cref{ch:coordination}), which is currently stated only as a
conjecture in this monograph.

The monograph's own worked constructions supply natural first test
cases rather than requiring new ones to be invented: \cref{con:witness}
and its extension (\cref{prop:extension}) give an exact, small
instance against which a simulator's discard behavior could be checked
directly, \cref{con:pipeline-instance} gives a named, domain-flavored
instance of the same size, and \cref{con:two-branch} gives the minimal
fan-out case. A useful early milestone for any implementation would be
reproducing these four constructions' discard sets exactly, before
attempting the larger, randomized relays needed to explore
\cref{cor:no-free-lunch}'s growth-rate question empirically.
